\chapter{So sánh và thảo luận}

\section{Cải tiến có khái quát qua ba bộ dữ liệu không?}

Để kiểm tra robustness, Bảng~\ref{tab:robustness-e4} so sánh cùng một cấu hình kết hợp E4 với baseline. Letter có $\alpha=0$ nên E4 tương đương HS trên cây chưa prune; Digits chọn $\lambda=0$ sau CCP; Covertype dùng cả $\alpha>0$ và $\lambda=10$.

\begin{table}[H]
    \centering
    \small
    \caption[Tác động của E4 trên ba bộ dữ liệu]{Tác động của E4 so với baseline trên ba chế độ dữ liệu.}
    \label{tab:robustness-e4}
    \setlength{\tabcolsep}{3.5pt}
    \begin{tabularx}{\textwidth}{@{}>{\raggedright\arraybackslash}p{2.15cm}>{\centering\arraybackslash}p{2.15cm}>{\centering\arraybackslash}p{1.75cm}>{\centering\arraybackslash}p{1.65cm}>{\raggedright\arraybackslash}X@{}}
        \toprule
        \textbf{Dataset} & \textbf{$\Delta$ Macro-F1} & \textbf{$\Delta$ Leaves} & \textbf{E4 Error} & \textbf{Diễn giải} \\
        \midrule
        Letter & $+0.11$ pp & $0.0\%$ & 13.12\% & HS cải thiện rất nhẹ; CV không chọn pruning \\
        Digits & $-0.71$ pp & $-54.9\%$ & 18.33\% & Cây nhỏ hơn nhưng mất điểm; mẫu nhỏ làm estimate kém ổn định \\
        \textbf{Covertype} & \textbf{$+0.19$ pp} & \textbf{$-31.7\%$} & \textbf{6.07\%} & \textbf{Trade-off rõ nhất giữa score, gap và complexity} \\
        \bottomrule
    \end{tabularx}
\end{table}

Kết quả không ủng hộ phát biểu ``HS luôn tăng score''. Lợi ích regularization phụ thuộc số mẫu, độ sâu cây và cấu trúc dữ liệu. Covertype có đủ dữ liệu và một cây baseline đủ lớn để CCP + HS tạo thay đổi đáng tin cậy về complexity; trên Digits, pruning mạnh làm mất chi tiết hữu ích.

\section{Benchmark mở rộng: Decision Tree so với ba họ mô hình}

Decision Tree là mô hình nghiên cứu chính. Random Forest, SVM--RBF và KNN chỉ là reference baselines đại diện cho ensemble tree, kernel và instance-based learning \cite{breiman2001rf,cortes1995svm,cover1967knn}. Bảng~\ref{tab:best-models} chỉ tóm tắt mô hình tốt nhất trên từng dataset; bảng đầy đủ nằm ở Phụ lục~\ref{app:full-benchmark}.

\begin{table}[H]
    \centering
    \small
    \caption[Mô hình tốt nhất trên từng bộ dữ liệu]{Mô hình tốt nhất theo test accuracy trong các cấu hình đã khảo sát.}
    \label{tab:best-models}
    \begin{tabular}{@{}llrr@{}}
        \toprule
        \textbf{Dataset} & \textbf{Mô hình tốt nhất} & \textbf{Accuracy} & \textbf{Macro-F1} \\
        \midrule
        Letter Recognition & Random Forest & 95.13\% & 95.12\% \\
        Handwritten Digits & SVM--RBF & 97.22\% & 97.18\% \\
        Covertype & Decision Tree & 93.89\% & 90.34\% \\
        \bottomrule
    \end{tabular}
\end{table}

Random Forest giảm phương sai hiệu quả trên Letter; SVM--RBF phù hợp với dữ liệu ảnh nhỏ và nhiều chiều; Decision Tree khai thác tốt dữ liệu bảng Covertype trong cấu hình hiện tại. Kết quả Random Forest Covertype 79.57\% thuộc implementation cuML và một cấu hình không được tìm kiếm siêu tham số toàn diện. Vì vậy, benchmark chỉ mô tả pipeline đã thử nghiệm và không hàm ý Decision Tree luôn vượt Random Forest hoặc SVM.

\section{Tính tái lập và công bằng thực nghiệm}

Các biện pháp kiểm soát gồm:

\begin{itemize}
    \item một phép chia 80/20 stratified và \texttt{random\_state=42} dùng chung;
    \item preprocessing chỉ fit trên train, test không dùng để chọn tham số;
    \item $\alpha$ và $\lambda$ được chọn bằng 3-fold CV trên toàn bộ train;
    \item cùng test set được dùng cho bốn mô hình trên từng dataset;
    \item baseline Covertype được kiểm tra thêm bằng 5-fold CV;
    \item phiên bản thư viện, backend và hardware được lưu cùng artifact.
\end{itemize}

5-fold CV của E0 Covertype cho accuracy $93.30\%\pm0.11$ điểm phần trăm và macro-F1 $89.23\%\pm0.11$ điểm phần trăm. Độ lệch chuẩn nhỏ cho thấy baseline tương đối ổn định giữa các fold. Ngược lại, một cấu hình đặt trước với \texttt{max\_depth=20} và \texttt{min\_samples\_leaf=5} chỉ đạt macro-F1 $83.88\%\pm0.28$, xác nhận regularization phải được chọn bằng validation thay vì áp dụng tùy ý.

\section{Môi trường tính toán và diễn giải runtime}

Decision Tree và HS dùng scikit-learn trên CPU \cite{pedregosa2011sklearn}. Random Forest, SVM và KNN trong notebook benchmark dùng RAPIDS cuML trên GPU \cite{rapidsCuml}. Phiên benchmark được cấu hình Kaggle GPU x2 và metadata cuML ghi nhận Tesla T4; các notebook Decision Tree ghi nhận Tesla P100 khả dụng nhưng compute device vẫn là CPU.

Do backend, implementation và thuật toán khác nhau, runtime chỉ phản ánh chi phí thực tế của từng pipeline. Nghiên cứu không diễn giải các số này như speedup CPU--GPU thuần túy và không khẳng định hai GPU được sử dụng song song.

\section{Threats to validity}

\begin{enumerate}
    \item \textbf{Một holdout chính:} ngoài 5-fold CV của baseline Covertype, kết quả test chính dùng một seed; mức tăng nhỏ cần được lặp lại với nhiều seed.
    \item \textbf{Tuning chưa exhaustive:} candidate grid tập trung vào Decision Tree improvements; RF, SVM và KNN chưa được tối ưu toàn diện.
    \item \textbf{Backend không đồng nhất:} Decision Tree chạy CPU còn ba mô hình đối chứng chạy GPU, nên runtime không phải benchmark phần cứng có kiểm soát.
    \item \textbf{Effect size nhỏ:} chênh lệch score giữa E0, E2, E3 và E4 nhỏ; chưa có kiểm định thống kê nhiều lần lấy mẫu để khẳng định ưu thế tuyệt đối.
\end{enumerate}

Những giới hạn này thu hẹp phạm vi kết luận nhưng không làm mất kết quả chính về complexity: E4 giảm 31.7\% số lá trên cùng test protocol mà vẫn giữ macro-F1 gần như không đổi và giảm gap.
